\section{Categorification of integrable highest weight $\U({\widehat{\gl}_p})$-modules}
\label{sec:Cat}

\subsection{Higher level Fock spaces}

Denote by $\Par$ the set of all partitions. Given $\ell\ge 1$ and $s\in \Z$, let 
\begin{align} \label{eq:Zls}
\Z^\ell(s) =\big\{\bfs =(s_1, \ldots, s_\ell) \in \Z^\ell \mid s_1+ \ldots+ s_\ell =s \big\}.
\end{align}
One has a Fock space $\cF^{s+\infty}$ of level $\ell, p$ and charge $s$, which is equipped with a natural basis $\{|\bfla, \bfs \rangle \}$, where $\bfla =(\lambda^{(1)}, \ldots, \lambda^{(\ell)}) \in \Par^\ell$ and $\bfs \in \Z^\ell(s)$. It admits an action of $\U_v(\widehat\gl_p)$ from \eqref{eq:affine_glp} (and also a commuting action of $\U_v(\widehat\sll_\ell)$ which will be ignored here), cf. \cite{JMMO91, Ugl00}. 

We have a decomposition as $\U(\widehat\gl_p)$-modules
\[
\cF^{s+\infty} = \bigoplus_{\bfs \in \Z^\ell(s)} \cF_\bfs, 
\]
where $\cF_\bfs$ is spanned by $|\bfla, \bfs \rangle$, for all $\bfla \in \Par^\ell$. 

Uglov \cite{Ugl00} defined a bar involution on $\cF_\bfs$, for each $\bfs$ (and hence on $\cF^{s+\infty}$) such that
$\overline{v}=v^{-1}, \overline{|\emptyset, \bfs \rangle} =|\emptyset, \bfs \rangle$, $\overline{ux} =\overline{u}\,\overline{x}$, and $\overline{B_{-r}x} =\overline{B_{-r}}\overline{x}$, for all $x\in \cF_\bfs$, $u \in \U_v(\widehat\sll_p)$ and $r>0$. 

\begin{theorem} \cite{Ugl00}
Let $\bfs\in \Z^\ell(s)$ for $s\in \Z$. 
Then there exists a canonical basis $\bfB_{\bfs} =\{b_{\bfla,\bfs} \mid \bfla \in \Par^\ell \}$ 
%and a dual canonical basis $\bfB_{\bfs}^- =\{b_{\bfla,\bfs}^- \mid \bfla \in \Par^\ell \}$ 
on $\cF_\bfs$ which are characterized by 
\begin{enumerate}
    \item 
$\overline{b_{\bfla,\bfs}} =b_{\bfla,\bfs}$,
    \item
$b_{\bfla,\bfs} \in |\bfla, \bfs \rangle + \sum\limits_{\bfmu \in \Par^\ell,\bfmu \lhd \bfla} v\Z[v] |\bfmu, \bfs \rangle$.
% , where $\lhd$ denotes the dominance order on multi-partitions.  
%\qquad $b_{\bfla,\bfs}^- \in |\bfla, \bfs \rangle + \sum\limits_{\bfmu \in \Par^\ell} v^{-1} \Z[v^{-1}] |\bfmu, \bfs \rangle$.
\end{enumerate}
\end{theorem}

Let $\omega_k$ be the $k$th fundamental weight, and denote $\omega_\bfs =\omega_{s_1} +\ldots +\omega_{s_\ell}$.
Denote by $\Lambda (\omega_\bfs)$ the integrable $\U(\widehat{\mathfrak {gl}}_p)$-module of highest weight $\omega_\bfs$ with highest weight vector $v^+_{\omega_\bfs}$. Denote by $\Lambda (\omega_\bfs)\!_\cA :=\UA^-(\widehat{\mathfrak {gl}}_p)\,  v^+_{\omega_\bfs}$ the $\cA$-form and by $\Lambda (\omega_\bfs)_\Z$ the specialization at $v=1$. 
Recall the canonical basis $\bfB$ on the generic Hall algebra $\Hall$ and the $\U^-_\cA(\widehat{\mathfrak {gl}}_p)$-module surjection $\varpi_\bfs: \Hall \rightarrow \Lambda(\omega_\bfs)\!_\cA$ thanks to the identification $\U^-_\cA(\widehat\gl_p) \cong \Hall$; see \eqref{eq:isom}.

Denote by $\mathcal C (\omega_\bfs)$ the subgraph for $\Lambda (\omega_\bfs)$ of the crystal graph for $\cF_\bfs$ with vertex set in $\Par^\ell$; cf. \cite{Kas91, Sch00}. 

\begin{theorem} \cite[Theorem 4.2]{Sch00}
\label{thm:Sch00}
    We have $\bfB \cdot |\emptyset, \bfs \rangle \subset \bfB_{\bfs} \cup \{0\}.$ Moreover, for $\bfla\in \mathcal C (\omega_\bfs)$, we have $b_{\bfla,\bfs} \in |\bfla, \bfs \rangle + \sum\limits_{\bfmu \in \Par^\ell, \bfmu\lhd\bfla} v\N[v] |\bfmu, \bfs \rangle$.
\end{theorem}
Thus $\bfB \cdot |\emptyset, \bfs \rangle \big\backslash \{0\}$ forms a $\C(v)$-basis for $\Lambda (\omega_\bfs)$ and an $\cA$-basis for $\Lambda (\omega_\bfs)_\cA$ (called a canonical basis). For $\bfla\in \mathcal C (\omega_\bfs)$, we write
\begin{align} \label{eq:decomposition}
    b_{\bfla,\bfs} = \sum_{\bfmu \unlhd \bfla} d_{\bfmu,\bfla}(v) |\bfmu, \bfs \rangle,
\end{align}
where $d_{\bfla,\bfla}(v)=1$, and $d_{\bfmu,\bfla}(v) \in v\N[v]$, for $\bfmu\lhd\bfla$. Also set $d_{\bfmu,\bfla}(v)=0$, for $\bfmu\ntrianglelefteq\bfla$. We refer to $\bold{D}(v) =\big[d_{\bfu,\bfla}(v)\big]_{\bfmu,\bfla}$ or $\bold{D}_n(v) =\big[d_{\bfmu,\bfla}(v)\big]_{|\bfmu|=\bfla|=n}$ as the $v$-decomposition matrix for $\qW_{\eps^{2\bfs}}$ or for $\wS_{\eps^{2\bfs}}(n)$. 




\subsection{Cyclotomic categorification}

We continue to use the notation $\bfs \in \Z^\ell$ from \eqref{eq:Zls}, which differs from the notation used in $\eps^{2\bfs}$ (with $\bfs \in \I^\ell$) in earlier sections, but $\eps^{2\bfs}$ still makes sense. 
%only need $s_1\gg \cdots \gg s_\ell$ as in \cite{Yv06, Yv07} }
Thanks to Proposition \ref{prop:IndRes} we have the following. 
  
\begin{proposition} 
  \label{prop:lieaction}
There is a $\U_\Z(\widehat{\mathfrak {gl}}_p)$-module structure on $[\qW_{\eps^{2\bfs}}\modf]$ such that $e_{\bfi}$ and $f_{\bfi}$ act via  $[\iRes]$ and $[\iInd]$, respectively.
In particular, $[\Delta^\emptyset]$ is of weight $\omega_\bfs$.  This induces a $\U_\Z(\widehat{\mathfrak {gl}}_p)$-module structure on $[\qW_{\eps^{2\bfs}}\modf]^*$ such that %$e_{\bfi}$ and $f_{\bfi}$ act via 
$\phi\mapsto e_\bfi\phi, f_\bfi\phi$ with
\[ 
(e_\bfi \phi) ([M])=\phi ([\iInd(M)]), \qquad
(f_\bfi\phi) ([M])=\phi([\iRes(M)]). 
\]
\end{proposition}

\begin{proof}
    The second statement follows from the first one. Identifying $[\Delta^\lambda]$ with the Schur basis of the higher level Fock space $\cF_{\bfs,\Z}$ (which can be assumed to be a tensor product without loss of generality), the first statement follows by comparing  Proposition \ref{prop:IndRes} and the actions of $e_\bfi$ and $f_\bfi$ on the Fock space  and applying the double Hall construction; cf. \cite[\S 6.2]{VV99} and \cite{X97}.
    %(see also \cite[Lemma 7.9]{SW11}). 
\end{proof}


\begin{theorem}  \label{thm:cyc_cat}
\begin{enumerate}
    \item 
There exists an isomorphism of $\U_\Z(\widehat{\mathfrak {gl}}_p)$-modules
\[
\jmath_\bfs: \Lambda (\omega_\bfs)_\Z   \stackrel{\cong}{\longrightarrow}  [\qW_{\eps^{2\bfs}} \modf]^*, \qquad v^+_{\omega_\bfs}\mapsto[\Delta^\emptyset]^*.
\]
\item
We have the following commutative diagram of $\Uglzhalf$-module homomorphisms: 
\begin{equation}  
\label{CD:cat}
\begin{tikzcd}
\HallZ \ar[r,"\jmath_p"]\ar[d,two heads,"\varpi_\bfs"]&{[}\qAW \modf_\eps{]}^* \ar[d,two heads,"\pi_\bfs^*"]\\
\Lambda (\omega_\bfs)_\Z  \ar[r,"\jmath_\bfs"]&{[}\qW_{\eps^{2\bfs}} \modf{]}^* 
\end{tikzcd} 
\end{equation}
\item
 The isomorphism $\jmath_\bfs$ sends the canonical basis to the basis dual to the classes of simple modules. 
 \end{enumerate}
\end{theorem}

\begin{proof}
Let us first explain the maps in the following diagram:
\begin{equation}  
\label{CD:modules}
\begin{tikzcd}
\HallZ \ar[r,two heads,"\varpi_\bfs"]\ar[d,"\jmath_p"]& \Lambda (\omega_\bfs)_\Z \ar[r,hook] \ar[d, dashed,"\jmath_\bfs"] &\cF_{\bfs,\Z} \\
{[}\qAW \modf_\eps{]}^* \ar[r,two heads,"\pi_\bfs^*"]&  {[}\qW_{\eps^{2\bfs}} \modf{]}^*  \ar[ur,"\psi^*"]
\end{tikzcd} 
\end{equation}
The homomorphism $\jmath_p$ was given in Proposition \ref{prop:lieaction}, $\varpi_\bfs$ is the natural surjective homomorphism by the identification $\HallZ =\Uglzhalf$, $\pi_\bfs^*$ is the natural surjective homomorphism induced from the embedding $\pi_\bfs: [\qW_{\eps^{2\bfs}} \modf] \rightarrow [\qAW \modf_\eps]$, and $\Lambda (\omega_\bfs)_\Z \hookrightarrow\cF_\bfs$ is the natural $\U(\widehat{\mathfrak {gl}}_p)_\Z$-module embedding sending $v^+_{\omega_\bfs}\mapsto |\emptyset,\bfs\rangle$. Note that the $\U(\widehat{\mathfrak {gl}}_p)_\Z$-action on $\cF_\bfs$ is interpreted in \cite{DX17} as the double Hall algebra action, extending the identification of the $\U^-(\widehat{\mathfrak {gl}}_p)_\Z$-action with $\Hall$-action on $\cF_\bfs$.

We identify $\cF_{\bfs,\Z}$ with the Grothendieck group of the semisimple category $\qW_{t^\bullet\eps^{2\bfs}}\modf$ for a semisimple parameter of the form $t^\bullet\eps^{2\bfs}$ deforming $\eps^{2\bfs}$ (see \cite[(4.5)]{Ar96} and Proposition \ref{prop:semisimple}), and there is a natural surjection $\psi:\cF_\bfs \rightarrow [\qW_{\eps^{2\bfs}} \modf]$, which sends the Schur basis elements $|\bfla,\bfs\rangle$ (identified with Specht modules) to $[\Delta^{\bfla}]$, for $\bfla\in \Par^\ell$; this is a $\U(\widehat{\mathfrak {gl}}_p)$-module homomorphism as they commute with all $f_\bfi$ and $e_\bfi$ (arising from restriction and induction functors). Dualizing $\psi$ and noting that $\cF_{\bfs,\Z}$ is self-dual with orthnormal Schur basis, we obtain the injective homomorphism $\psi^*: [\qW_{\eps^{2\bfs}} \modf]_\C^*\rightarrow \cF_{\bfs,\C}$ sending $[\Delta^\emptyset]^* \mapsto |\emptyset,\bfs\rangle$.

It remains to construct the isomorphism $\jmath_\bfs$ in order to complete the commutative diagram \eqref{CD:modules}. As a $\Uglzhalf$-module, $\HallZ$ is cyclic and so is ${[}\qW_{\eps^{2\bfs}} \modf{]}^*$ with $[\Delta^\emptyset]^*$ as a cyclic vector. Passing through a base change from $\Z$ to $\C$, $\cF_{\bfs,\C}$ is a semisimple $\U(\widehat{\mathfrak {gl}}_p)$-module, and
hence as its cyclic highest weight submodule $[\qW_{\eps^{2\bfs}} \modf]_\C^*$ must be a simple $\U(\widehat{\mathfrak {gl}}_p)$-module of highest weight $\omega_\bfs$. Therefore, we obtain an isomorphism of $\U(\widehat{\mathfrak {gl}}_p)$-modules $\jmath_\bfs: \Lambda (\omega_\bfs)   \rightarrow  [\qW_{\eps^{2\bfs}} \modf]^*_\C,$ $v^+_{\omega_\bfs} \mapsto[\Delta^\emptyset]^*$, completing the commutative diagram \eqref{CD:modules}. This proves Parts (1)--(2).

It remains to prove (3). The embedding $\pi_\bfs: [\qW_{\eps^{2\bfs}} \modf] \rightarrow [\qAW \modf_\eps]$ sends simples to simples. By Schiffman's Theorem \ref{thm:Sch00}, $\varpi_\bfs: \HallZ \rightarrow \Lambda(\omega_\bfs)_\Z$ sends the canonical basis to canonical basis or $0$. By Theorem \ref{thm:Cat_affine}, $\jmath_p: \HallZ \rightarrow [\qW_{\eps^{2\bfs}}\modf]^*$ sends the canonical basis to the basis dual to the classes of simple modules. Collecting all these together, we conclude that $\jmath_\bfs: \Lambda (\omega_\bfs)_\Z\rightarrow  [\qW_{\eps^{2\bfs}} \modf]^*$ sends the canonical basis to the basis dual to the simple modules.
\end{proof}

\begin{rem}
\begin{enumerate}
\item 
One expects that the same categorification results as Theorem \ref{thm:cyc_cat} for $q$ generic can be obtained for degenerate cyclotomic webs over $\C$ \cite{SW25web}. 
\item
One also expects that Parts (1)-(2) of Theorem \ref{thm:cyc_cat} remain valid while Part (3) on canonical basis will fail for module categories of degenerate affine and cyclotomic webs over a field of positive characteristic $p$. 
\end{enumerate}
\end{rem}

\begin{rem} \label{rem:Cartan_pairing}
    Denote by $\qW_{\eps^{2\bfs}} \prmod$ the category of finite-dimensional projective $\qW_{\eps^{2\bfs}}$-modules. Via the Cartan pairing, one can replace $[\qW_{\eps^{2\bfs}}\modf]^*$ and the duals of the simple modules in Theorem \ref{thm:cyc_cat}  by $[\qW_{\eps^{2\bfs}}\prmod]$ and the classes of indecomposable projective modules.
\end{rem}

The category $\qW_{\eps^{2\bfs}} \modf$  has enough projectives and every projective module admits a cell module filtration since $\qW_{\eps^{2\bfs}}$ is a cellular algebra. 

\begin{corollary}
\label{cor:PIM}
    For $\bfla \in \mathcal C(\omega_\bfs)$, there exists a unique (up to isomorphism) indecomposable projective $\qW_{\eps^{2\bfs}}$-module $P^{\bfla}$, whose sections are all of the form $\Delta^\bfmu$ for $\bfmu\unlhd \bfla$ and $\Delta^\bfla$ appears exactly once, such that  
\begin{align}
    [P^{\bfla}] = \sum_{\bfmu \unlhd \bfla} d_{\bfmu,\bfla}(1) \,[\Delta^{\bfmu}].
\end{align}
%Such $P^\bfla$ can be determined recursively by the dominance order on $\Par^\ell$.
\end{corollary}

\begin{proof}
    It follows by Theorem \ref{thm:cyc_cat}, in light of the $q$-decomposition matrix \eqref{eq:decomposition} and Remark~ \ref{rem:Cartan_pairing}(1). 
\end{proof}
Denote by $L^\bfla$ the simple head of $P^\bfla$. Then we have the surjective homomorphisms of $\qW_{\eps^{2\bfs}}$-modules 
\[
P^{\bfla} \twoheadrightarrow \Delta^{\bfla} \twoheadrightarrow L^{\bfla}.
\]
We have the following classification result. 
\begin{corollary} \label{cor:classifysimple}
    The $L^\bfla$, for $\bfla \in \mathcal C(\omega_\bfs)$, forms a complete list of pairwise non-isomorphic simple modules in $\qW_{\eps^{2\bfs}} \modf$. 
\end{corollary}

There is an idempotent functor which we refer to as a Schur functor, for each $n$, $1_{(1^n)}: \qW_{\eps^{2\bfs}}(n)\modf \rightarrow \HH_{\eps^{2\bfs}}(n)\modf$. Then  $1_\bfs: =\oplus_{n\in \N}1_{(1^n)}$ leads to a Schur functor (which is surjective) $1_\bfs: [ \qW_{\eps^{2\bfs}}\modf]
\rightarrow [\HH_{\eps^{2\bfs}}\modf]$
and hence an injective map 
$1_\bfs^*:[\HH_{\eps^{2\bfs}}\modf]^*\rightarrow [ \qW_{\eps^{2\bfs}}\modf]^*$. Moreover, for $i\in \I$, there are functors $i$-$\Res$ and $i$-$\Ind$ on $\HH_{\eps^{2\bfs}}\modf$ \cite{Ar96} defined in the same way as those for $ \qW_{\eps^{2\bfs}}\modf$. Ariki's categorification theorem \cite[Theorem 4.4]{Ar96} shows that these functors define an $\widehat{\mathfrak{sl}_p}$-module structure on
$[\HH_{\eps^{2\bfs}}\modf]^*$ such that $[\HH_{\eps^{2\bfs}}\modf]^*\cong V(\omega_\bfs)_\Z$, where $V(\omega_\bfs)$ denotes the integrable  $\widehat{\mathfrak{sl}}_p$-module of highest weight $\omega_\bfs$. 

\begin{proposition} 
\label{prop:SchurHecke}
We have a commutative diagram of $\U_\Z(\widehat{\mathfrak{sl}}_p)$-module homomorphisms:
\begin{equation}   \label{CD:HeckeWeb}
\begin{tikzcd}
V (\omega_\bfs)_\Z \ar[d,hook,"\iota_\bfs"]\ar[r,"\phi_\bfs"] &{[}\HH_{\eps^{2\bfs}}\modf {]}^*\ar[d,hook,"1_\bfs^*"] \\
\Lambda (\omega_\bfs)_\Z  \ar[r,"\jmath_\bfs"]& {[}\qW_{\eps^{2\bfs}} \modf{]}^* 
\end{tikzcd} 
\end{equation}
where $\phi_\bfs$ is Ariki's isomorphism \cite{Ar96}. Moreover, the embedding $\iota_\bfs$ preserves the canonical basis while $1_\bfs^*$ preserves the basis dual to the simple modules. 
\end{proposition}

\begin{proof}
Since $1_\bfs \qW_{\eps^{2\bfs}} 1_\bfs= \HH_{\eps^{2\bfs}}$ \cite{SSW25}, 
  by definition we have  
  $i\text{-Res}1_\bfs (M)= 1_\bfs i\text{-Res}(M)$ and $i\text{-Ind}1_\bfs (M)= 1_\bfs i\text{-Ind}(M)$, for any $M\in \qW_{\eps^{2\bfs}} \modf$.
  This shows that $1_\bfs^*$ is a homoporphism of $\widehat{\mathfrak{sl}_p}$-modules. 
  Then by comparing the actions of $i$-Res and $i$-Ind on the cell modules $ S^{\lambda}=1_\bfs \Delta^\lambda$ in \cite[Lemma 2.1]{Ar96} and that on $\Delta^\lambda $ in Proposition \ref{prop:IndRes}, we see that  the diagram is commutative. The rest follows from Theorem \ref{thm:cyc_cat}.
\end{proof}

\begin{rem}
Proposition \ref{prop:SchurHecke} shows  that the $v$-decomposition matrix for a cyclotomic $q$-W-Schur algebra contains the counterpart for a cyclotomic Hecke algebra as a submatrix. At level one, this was due to \cite{LT96, VV99}. 
\end{rem}

\begin{rem}
The canonical basis on $\Lambda(\omega_\bfs)$ and thus the $v$-decomposition matrices depend only on the parameters $\bfs$ modulo $p$. This differs from the observation in \cite{Yv06} on the canonical basis of Uglov's higher Fock spaces.
\end{rem}


\subsection{Some open problems}
\label{ssec:open}

There has been much further advances on interaction among Hecke algebras, quantum groups and categorification since the works of Lascoux-Leclerc-Thibon and Ariki. We refer to the survey of Kleshchev \cite{Kle10} for extensive references. The cyclotomic $q$-Schur algebras have also been widely studied. The affine and cyclotomic $q$-web categories sit in between these well-studied algebras as indicated by the diagram \eqref{CD:categories}. It is reasonable to ask for various W-Schur analogues of the topics in the other two contexts, and we list some of them.

\begin{enumerate}
    \item Construct a Jantzen filtration for the cell modules of cyclotomic $q$-W-Schur algebras and compute the Jantzen sum formula. 
    We expect that the Jantzen filtration is compatible with the $v$-decomposition matrix for cyclotomic $q$-W-Schur algebras. 
    \item Classify the blocks for cyclotomic $q$-W-Schur algebras. 
    \item Understand the relation between the crystal basis and modular branching rules for cyclotomic $q$-W-Schur algebras. The rich combinatorics of crystals often corresponds to different ways of parameterizing the simple modules. 
    \item Formulate a $\Z$-graded version.
    \item Realize an action of $\U_\Z^-(\widehat{\mathfrak{gl}}_p)$ on $\K^*$ geometrically so that 
    \eqref{eq:ijk} becomes a commutative diagram of $\U_\Z^-(\widehat{\mathfrak{gl}}_p)$-module isomorphisms. 
    \item  (This is a variant and an enhancement of (5).)  Construct algebraic structures on ${[}\qAW\modf_\eps{]}^*$ diagrammatically and on $\K^*$ geometrically so that 
    \eqref{eq:ijk} becomes a commutative diagram of algebra isomorphisms.
\end{enumerate}
\noindent {\bf Conflict of interest.}

On behalf of all authors, the corresponding author states that there is no conflict of interest.
The authors declare that the data supporting the findings of this study
are available within the paper.